\section{Related work}
\label{sec:related}

\textbf{Group-relative policy optimisation.} GRPO \citep{shao2024deepseekmath} replaces the learned
value function of PPO \citep{schulman2017ppo, ouyang2022instructgpt} with a baseline computed
across a group of sampled completions for the same prompt, and is now widely used for RL on
language models \citep{guo2025deepseekr1}. For dialogue this is attractive, since no value
function over conversation states is needed, but the entire learning signal is then the
\emph{relative} ranking of siblings, so the reward determines everything. We change what the
reward observes and nothing else.

\textbf{Multi-turn RL for dialogue.} That a dialogue turn should be credited for its downstream
consequences is as old as RL-based dialogue management \citep{levin2000stochastic,
singh2002optimizing}; it was brought to neural dialogue generation by \citet{li2016deeprl} and
pursued since with utterance-level value functions \citep{zhou2024archer}, offline RL on
LLM-imagined conversations \citep{hong2023imagined}, self-reinforcement on role-played
interactions filtered by an LLM-rated social goal \citep{wang2024sotopiapi}, preference feedback
over whole conversations \citep{shani2024multiturn}, regression of relative future rewards from
on-policy continuations \citep{gao2024refuel}, and step-wise critics given training-time
information \citep{zhou2025sweetrl}. GRPO itself has been carried into multi-turn settings by
turn-level credit assignment in tool-use agents \citep{wei2025multiturn} and by rollouts against
LLM-simulated users \citep{qian2025userrl}; both redistribute the environment's reward over
turns. Look-ahead instead changes what the reward is computed on: no value function, no new
loss, a Monte Carlo estimate of a turn's downstream value from one rollout under the current
policy, as VinePPO does for reasoning chains \citep{kazemnejad2024vineppo}.

\textbf{Look-ahead and search in preference learning.} Simulated continuations drive
inference-time dialogue planning, as prompt-based MCTS \citep{yu2023promptmcts} or planning in a
learned semantic space \citep{chen2025broaden}, and enter preference learning as MCTS-guided
search \citep{xie2024mctsdpo}; related mechanisms score partial trajectories with process reward
models \citep{lightman2024letsverify} or keep the best and worst of $N$ samples by a reward model
\citep{pace2024westofn}. The closest prior work is PTO \citep{baruch2025pto}, which introduced
$K$-turn look-ahead for MI dialogue: candidate turns branched at each patient turn, each scored by
an oracle shown its simulated continuation, the best and worst becoming a DPO pair
\citep{rafailov2023dpo}. We keep the lever and drop the tree and the pair selection.

\textbf{Reward hacking and LLM judges.} That optimising a proxy degrades the true objective is a
long-standing concern \citep{amodei2016concrete, skalse2022defining, pan2022effects, gao2023scaling},
and LLM judges \citep{zheng2023judging} and preference-trained reward models add distortions of
their own: they can reward sycophantic answers \citep{sharma2024sycophancy, perez2023discovering},
favour their own generations \citep{panickssery2024selfpreference}, and favour longer outputs
\citep{singhal2024long, dubois2024lengthcontrolled}. The oracle saturation of \S\ref{sec:measurement} is the
variance side of reward-model over-optimisation
\citep{gao2023scaling} on a bounded judge scale: a trained-against grader driven to its ceiling
loses discrimination over conversations that an independent grader with headroom still tells apart;
our second grader plays in evaluation the role reward-model ensembles play in training
\citep{coste2024ensembles}.

\textbf{Motivational interviewing.} MI \citep{miller1991mi} has established coding schemes
\citep{moyers2016miti}, validated session instruments \citep{hatcher2006waisr, larsen1979csq}, and
an expert-annotated corpus of high- and low-quality sessions \citep{wu2022annomi}. Prior work assesses
counselling quality automatically \citep{perezrosas2019goodcounselor, yosef2024assessing}, trains
counsellors against LLM-simulated patients \citep{steenstra2025scaffolding, wang2024patientpsi},
and codes the behaviour of LLM
therapists, which resemble low-quality therapy in advising where a counsellor would reflect
\citep{chiu2024bolt}; a smaller version of the same directive pattern appears in our look-ahead
policy (\S\ref{sec:behaviour}).
